This thesis estimates the probability that Argentina's public debt fails
to stabilize in 2026-2035 and assesses the role of growth, the real
exchange rate, country risk, the primary balance and United States
monetary policy. A structural vector autoregression with exogenous
variables (SVARX) is estimated with quarterly data for 1996-2025 on
gross federal government debt, the perimeter of the International
Monetary Fund (IMF) debt sustainability analysis, with the Federal
Reserve (Fed) policy rate and global financial volatility as exogenous
variables. Debt is simulated around the IMF's 2026 baseline scenario,
and non-stabilization is measured with the institution's operational
definition and, as a complement, with an annual metric that includes the
valuation of foreign-currency debt in each year. Debt consolidated with
the Central Bank is an extension.

With the IMF definition, the probability is 3.1\% in 2031 (90\%
interval: 0.0\% to 6.3\%) and 5.1\% in 2035 (0.0\% to 10.2\%), against
about 1\% and less than 3\% reported by the IMF; the intervals include
those figures and the hypothesis that the probability is higher is not
supported. With the annual metric it is 27.7\% in 2031 (20.6\% to
34.7\%) and 29.0\% in 2035 (25.0\% to 36.1\%). Non-fiscal shocks
contribute 96.8 percentage points more than the fiscal shock to the
variance of 2035 debt. Repeating the 2022-2025 path of the Fed rate
raises the 2035 probability by 18.4 points, and the hypothesis on the
Fed is partially supported: the effect on expected debt includes zero
and 77\% of the increase comes from the direct cost of
dollar-denominated debt.

On the IMF metric, the thesis does not find that the institution
underestimates the risk; the difference appears once the annual
valuation of foreign-currency debt is included, which that metric
excludes. With either measure, most of the uncertainty comes from
external, exchange-rate and risk shocks.

Keywords: public debt, fiscal sustainability, SVARX, fan chart,
sovereign risk, Argentina.

JEL classification: C32, E62, F34, H63, H68.
